\section{模型假设}
% 篇幅上限：一页；条数 **3--5 条（最多 5 条）**，每条 1.5--2 行（约 55--80 字）。
% 假设节只收「建模约定」。题面已给定的条件、物理/数学必然（如「高度>0 才能起爆」）
%    不放进本节 —— 那是题面重述，放了要判返修。判据口径移到正文判据定义处。
%    每条假设末尾用一句话说明「放宽它会影响哪一问」。
%
% 排面（照参考件 `框架模板.pdf` 第 3 页「三、模型假设」的样子）：
%   ① 标题下先写一句引入语「为了适当地对模型进行合理简化，本文给出如下假设：」，
%      这句顶格（用 \noindent，不留两格缩进）。
%   ② 条目编号用 「假设 N：」，**不要项目符号**（不要 `•`）——
%      `[label=假设 \arabic*：]`；条目标签缩进两格、折行顶格（与正文段落同款）。
%   ③ 条目之间不加空距：与正文段落之间的距离一致（全局 `\setlist` 已把 itemsep/parsep 清零）。
%   参考件量值：引入句 x0=94.9、条目 x0=102.7、折行 x0=88.9、条目间距 19.8pt（= 正文行距）。

\noindent 为了适当地对模型进行合理简化，本文给出如下假设：

\begin{enumerate}[label=\textbf{假设\chinese{enumi}：}]
  \item <建模约定一>。放宽影响：<哪一问、怎么变>。
  \item <建模约定二>。放宽影响：<…>。
  \item <建模约定三>。放宽影响：<…>。
\end{enumerate}
